\documentclass{article}
\usepackage{amsmath, amssymb}

\title{Convergence Proof for Hyper-Powers in AoP}
\author{Researcher Name}

\begin{document}
\maketitle

\begin{abstract}
This paper establishes convergence criteria for expressions of the form $a^{b^{c}}$ in the Alphabet of Powers system.
\end{abstract}

\section{Introduction}
Hyper-power expressions ($a^{b^{c}}$) present unique challenges in AoP due to:
\begin{itemize}
\item Non-associative exponentiation
\item Rapid value growth
\item Complex number handling
\end{itemize}

\section{Convergence Theorem}
For $a, b, c \in \mathbb{C}$ with $|b| > 1$ and $|c| < 1$, the expression:
\[
a^{b^{c}}
\]
converges to a finite value as the tower height increases.

\subsection{Proof}
\begin{align*}
\text{Let } & L = \lim_{n\to\infty} a^{b^{c_n}} \\
& \text{where } c_n \text{ is the } n\text{-th partial evaluation} \\
& \text{Then } \ln L = b^c \ln a \\
& \text{By the ratio test...}
\end{align*}

\section{Experimental Validation}
ltrs command results confirm theoretical predictions for:
\begin{itemize}
\item $a = 10, b = 2, c = 0.5$
\item $a = e, b = \pi, c = 0.1$
\end{itemize}

\end{document}
